\documentclass[11pt]{article}
\usepackage[margin=1in]{geometry}
\usepackage{amsmath,amssymb}
\usepackage{booktabs}
\usepackage{graphicx}
\usepackage{xcolor}
\usepackage{listings}
\usepackage{url}
\usepackage[hidelinks]{hyperref}
\usepackage{longtable}

\lstset{basicstyle=\ttfamily\small,breaklines=true,columns=fullflexible,
        frame=single,framesep=4pt,xleftmargin=6pt,xrightmargin=6pt}

\newcommand{\sfv}{\textsf{sparkle-fv}}
\newcommand{\tool}[1]{\textsf{#1}}

\title{\sfv{}: An Agentic Harness for Autoformalization from
SystemVerilog to Sparkle HDL and AI-Assisted Formal Verification}
\author{The Sparkle HDL Project}
\date{July 2026}

\begin{document}
\maketitle

\begin{abstract}
We present \sfv{}, an agentic harness that (i)~\emph{autoformalizes}
SystemVerilog RTL into Sparkle HDL --- a hardware description language
embedded in the Lean~4 theorem prover --- and (ii)~performs block-level and
SoC-level formal verification with a portfolio of provers assisted by a
large language model. Following the ``everyday provers'' paradigm, every
verdict ships with an independently checkable artifact: safety proofs carry
$k$-induction certificates or explicit inductive invariants re-checked by an
SMT solver; bug reports carry counterexample traces that the harness
converts into stimulus and \emph{replays through Verilator} against the
original RTL. The AI layer (GLM~5.2 via OpenRouter) proposes candidate
invariants, repairs frontend incompatibilities, and synthesizes SoC-level
properties, but sits entirely outside the trusted base: all of its output is
filtered through a Houdini fixpoint and final Z3 re-checks, so it can
accelerate proofs but cannot produce a wrong ``safe.'' Applied to a
Linux-booting RV32IMA SoC (136 registers, 1{,}878 state bits, 6 memories),
the harness proved all eight SoC-level architectural invariants unboundedly
in about one minute and generated a design-verification convergence dossier
with 100\% cone-of-influence coverage; property synthesis further exposed
\emph{two previously unknown bugs} in the SoC (a \texttt{mstatus.MPP} WARL
violation reaching a reserved privilege mode, and a missing
instruction-address-misaligned trap). On a 13-design HWMCC-style benchmark
suite with 19 seeded bugs, \sfv{} matches the solve count of the
\tool{yosys-smtbmc} baseline (the engine behind SymbiYosys) under identical
elaborations while proving the safe SoC variant $34\times$ faster, and ---
uniquely --- confirms all 19 found bugs in Verilator simulation.
\end{abstract}

\section{Introduction}

Formal verification of RTL has long been split between two worlds:
push-button model checkers that scale but return opaque verdicts, and
interactive theorem provers that produce durable, auditable proofs but
demand months of expert effort. Recent progress in AI-assisted theorem
proving suggests a third path --- Pramaana Labs' essay \emph{The Age of
Everyday Provers}~\cite{pramaana} articulates it as a loop: translate domain
artifacts into formal representations, search the solution space with
provers and solvers, and return \emph{proof artifacts that domain experts
can inspect}, with AI accelerating every step but vouching for none of them.

\sfv{} instantiates this loop for hardware. Its inputs are ordinary
SystemVerilog designs; its outputs are (a)~Sparkle HDL models in Lean~4 with
machine-checkable proof obligations, (b)~unbounded safety certificates or
Verilator-confirmed counterexamples, and (c)~a \emph{convergence dossier}
--- the evidence package a design-verification (DV) team needs to sign off
a block or an SoC.

\paragraph{Contributions.}
\begin{enumerate}
  \item An \textbf{autoformalization pipeline} from SystemVerilog (via Yosys
        elaboration) to Sparkle HDL Lean~4 sources, covering the full
        word-level RTL cell library including memories and the complete
        flip-flop family, with proof-obligation skeletons that follow the
        Sparkle verification framework's proof patterns
        (\S\ref{sec:autoformal}).
  \item A \textbf{sound-by-construction AI proving architecture}: BMC,
        $k$-induction, and Houdini invariant synthesis over the
        HWMCC-standard BTOR2 format, where LLM-proposed lemmas and
        properties are admitted only after passing the same Z3 checks as
        deterministic ones (\S\ref{sec:engine}).
  \item \textbf{Bug surfacing with simulation confirmation}: counterexample
        traces are compiled into SystemVerilog testbenches and replayed
        through Verilator against the original RTL; a finding is reported
        as confirmed only when the assertion fires in simulation
        (\S\ref{sec:cex}).
  \item \textbf{Full-SoC convergence}: automatic property synthesis
        (structural mining plus LLM proposals), per-property dispositions
        with explicit proof radii, cone-of-influence coverage with named
        gaps, and an assumption audit, rendered as a reviewable dossier
        (\S\ref{sec:convergence}).
  \item An \textbf{empirical evaluation} on a 13-benchmark suite (12 block
        designs + a Linux-booting RV32IMA SoC, 19 seeded bugs with known
        minimal counterexample depths) against \tool{yosys-smtbmc} on
        identical elaborations, plus two real SoC bugs found during
        property authoring (\S\ref{sec:eval}).
\end{enumerate}

\section{Background}

\paragraph{Sparkle HDL.} Sparkle is a hardware description language embedded
in Lean~4: designs are pure functional programs over a \texttt{Signal}
abstraction (streams indexed by clock cycles), compiled through a word-level
IR to synthesizable SystemVerilog, C++ simulation, or a JIT backend. Because
the host language is a theorem prover, functional correctness theorems live
next to the design; the repository carries 60+ proved theorems over designs
ranging from arbiters to a BitNet inference ASIC, and its RV32IMA SoC boots
Linux under Verilator. \sfv{} gives this ecosystem an \emph{entry ramp}:
existing SystemVerilog can be lifted into Sparkle's IR and proof framework
mechanically.

\paragraph{BTOR2 and word-level model checking.} BTOR2 is the word-level
transition-system format of the Hardware Model Checking Competition (HWMCC).
Yosys emits it from elaborated RTL, mapping immediate assertions to
\texttt{bad} states and assumptions to \texttt{constraint} nodes. \sfv{}
adopts BTOR2 as its proving interface: the harness therefore runs unchanged
on any HWMCC-format model, and its Yosys frontend recipe pins one shared
semantic interpretation for the prover, the baseline tool, and the
autoformalizer.

\paragraph{Proof methods.} Bounded model checking (BMC) unrolls the
transition relation $k$ steps and asks a SAT/SMT solver for a path to a bad
state --- complete for bug finding up to the bound. $k$-induction
strengthens simple induction with $k$ consecutive safe frames and
(optionally) simple-path constraints. Houdini~\cite{houdini} computes the
\emph{greatest inductive subset} of a candidate-lemma set in a linear number
of solver calls; seeding the negated properties into the candidate set turns
a surviving fixpoint into an inductive invariant that implies safety.

\section{Architecture}\label{sec:arch}

\begin{figure}[h]
\begin{lstlisting}
                        agentic loops (repair / prove-refine / converge)
                        ================================================
SystemVerilog -> desugar -> Yosys --+--> BTOR2 --> Z3 transition system
   (frontend)   (+LLM repair)       |               |-- BMC (incremental)
                                    |               |-- k-induction
                                    |               +-- Houdini <-- lemmas:
                                    |                     templates + GLM-5.2
                                    +--> JSON netlist --> Sparkle HDL (Lean 4)
                                    |                     + proof obligations
                                    +--> SMT2 --> yosys-smtbmc (baseline)
BUG verdict -> trace.json + tb_replay.sv -> Verilator --assert -> CONFIRMED?
\end{lstlisting}
\caption{\sfv{} pipeline. All three elaboration outputs derive from one
Yosys recipe, so the prover, the baseline, and the autoformalizer see the
same design semantics.}
\label{fig:arch}
\end{figure}

The harness is a six-stage pipeline (Fig.~\ref{fig:arch}) in which every
stage consumes and produces a concrete artifact --- there is no hidden state,
so any stage can be re-run, audited, or swapped. A single Yosys recipe
(\texttt{prep; flatten; setundef -undriven -zero -init}) produces BTOR2 for
\sfv{}'s engine, SMT2 for the \tool{yosys-smtbmc} baseline, and a JSON
netlist for autoformalization. The \texttt{setundef -init} step gives every
flip-flop an explicit zero initial value, aligning the formal models with
Verilator's two-state simulation --- the property that makes counterexample
replay deterministic and the tool comparison apples-to-apples.

\paragraph{Agentic loops.} Three outer loops make the harness agentic, each
degrading gracefully to a deterministic strategy when the LLM is
unavailable: \emph{frontend repair} (desugar unsupported constructs; on
residual elaboration errors, ask the LLM for a minimal semantics-preserving
rewrite, validated by Yosys equivalence checking), \emph{prove--refine}
(escalate from BMC through $k$-induction to Houdini; on failure, ask the LLM
for design-specific candidate lemmas and re-run the fixpoint), and
\emph{converge} (synthesize, prove, and audit a SoC property set,
\S\ref{sec:convergence}).

\paragraph{Soundness architecture.} The trusted base is Yosys elaboration
plus Z3. The LLM can propose lemmas, properties, and rewrites; every
proposal passes through checks in the trusted base (Houdini fixpoint and
final re-check; property elaboration and the same proving pipeline; Yosys
\texttt{equiv\_make}). A hallucinated lemma wastes solver time; it cannot
produce an unsound verdict.

\section{Autoformalization to Sparkle HDL}\label{sec:autoformal}

Yosys netlists are bit-indexed while Sparkle's IR is word-level, so the
formalizer first runs a \emph{bit-vector net reconstruction} layer: adjacent
bit ids are regrouped into word nets (with priority given to ports and
public names), and cell connections that straddle nets become explicit
\texttt{slice}/\texttt{concat} expressions.

Two emission modes share this layer. \textbf{\texttt{circuitm} mode}
(robust, any netlist) rebuilds the design through Sparkle's \texttt{CircuitM}
builder monad --- covering arithmetic, logic, comparison and shift cells,
\texttt{\$mux}/\texttt{\$pmux}, the complete DFF family ($\pm$sync/async
reset, enables, polarities, folded with Yosys' priority semantics), and
\texttt{\$mem\_v2} memories. \textbf{\texttt{signal} mode} (idiomatic, small
blocks) emits applicative \texttt{Signal} DSL in the style of the
repository's examples and falls back to \texttt{circuitm} when a construct
does not fit.

For each property (including \texttt{\$assert} cells auto-extracted from the
netlist), the formalizer emits a self-contained proof-obligation file: pure
\texttt{Inputs}/\texttt{State}/\texttt{nextState} models, a
\texttt{Reachable} inductive, a decidable property predicate, a proof-plan
comment, and theorems in both the invariant pattern and Sparkle's temporal
(LTL) vocabulary. Inductive invariants discovered by the SMT engine are
injected into a marked \texttt{candidateInv} slot, so an unbounded Lean
proof can start exactly where the SMT engine stopped. Translation is
validated structurally (the IR mirrors the netlist one-to-one) and, for
agentic \emph{source} rewrites, by Yosys equivalence checking.

\section{The Proof Engine}\label{sec:engine}

The engine operates on a Z3 transition system parsed from BTOR2 (memories
stay word-level as SMT arrays --- no bit-blasting). Frames are materialized
lazily by substitution, so BMC, induction, and Houdini share one model.

\paragraph{BMC.} A single incremental solver; each depth adds transition
constraints, and each bad state is checked under \texttt{push}/\texttt{pop}
so learned clauses persist. On SAT, the model is projected onto per-frame
inputs plus the initial state to form a replayable trace.

\paragraph{$k$-induction.} The base case delegates to BMC; the step case
asserts $k$ consecutive safe frames (plus pairwise-distinct simple-path
constraints when no memories are present) and asks whether frame $k$ can be
bad. UNSAT yields verdict \textsc{safe} with certificate ``$k$-inductive.''

\paragraph{Houdini with AI lemmas.} Deterministic templates generate the
classic strengthenings (power-of-two and small-constant bounds, one-hot and
at-most-one-hot, per-bit constancy, capped pairwise orderings). The LLM path
sends the SystemVerilog source and the induction-failure context to GLM~5.2
and receives lemmas in a small JSON DSL
(\texttt{ule/eq/ne/implies/onehot/bit/in}) compiled to Z3 with width
coercion. Both sources feed the same fixpoint: candidates not implied by the
initial states are dropped; then, repeatedly, Z3 is asked for one transition
that breaks any surviving candidate, and everything falsified in that model
is dropped, until UNSAT. If the seeded negated properties survive, the
surviving conjunction is an inductive invariant proving safety, written out
for independent re-checking.

\paragraph{Verdicts.} \textsc{bug} (trace, then Verilator confirmation),
\textsc{safe} (certificate), or \textsc{unknown} (explicit bound and reason,
enabling escalation: deeper BMC, more lemmas, or a Lean obligation). The
portfolio never conflates a bounded result with a proof.

\section{Bug Surfacing and Verilator Stimulus}\label{sec:cex}

A BMC frame corresponds to one clock cycle. The generated testbench drives
non-clock inputs to frame $k$'s values in the low phase of cycle $k$
(stable before the $k$-th rising edge, matching BTOR2's
input-feeds-transition semantics) and lets Verilator's \texttt{--assert}
catch the immediate assertion. Because both sides zero-initialize state,
replay is deterministic: \texttt{confirmed = true} means the SMT-level
counterexample reproduces in simulation of the \emph{original} RTL --- the
strongest evidence a formal finding can carry into a simulation-centric
flow, and one that requires no formal expertise to consume.

\section{Full-SoC Convergence}\label{sec:convergence}

Sign-off requires more than ``we ran some properties'': the DV team needs a
defensible claim that the property set covers the design and that every
property has a disposition. \texttt{sparkle-fv converge} produces that claim.

\paragraph{Property synthesis.} Structural mining derives conjectures from
the elaborated model itself --- FSM value-set validity, counter bounds at
observed maxima, one-hot integrity --- sampled from shallow symbolic
exploration. In parallel, the LLM reads the RTL and proposes architectural
invariants as SystemVerilog assertion expressions, which are compiled by
instrumenting a copy of the source \emph{one property at a time} (a proposal
that fails elaboration is dropped without poisoning the batch). User
assertions are picked up as-is. All three origins go through identical
checking: synthesis proposes, the prover disposes.

\paragraph{Dispositions.} Every property ends in exactly one state:
\emph{proven} (unbounded certificate), \emph{bounded} (explicit proof
radius), \emph{falsified} (counterexample + Verilator replay),
\emph{vacuous} (tautology screen: holds in all states, reachable or not ---
zero coverage contribution), or \emph{discarded} (a refuted mined
conjecture --- not a design bug).

\paragraph{Coverage and audit.} The dossier computes cone-of-influence (COI)
coverage --- the fraction of state bits in the COI of at least one proven
property --- and lists uncovered registers by name as the DV to-do list. It
also lists every environment assumption in force, since proofs are only as
strong as their assumptions.

\paragraph{RV32IMA SoC case study.} Applied to the repository's
Linux-booting RV32IMA SoC (136 state elements, 1{,}878 state bits, 6
word-level memories, 10 free inputs; assertions cover MMU/PTW FSM validity,
pipeline control consistency, and AMO/CSR legality):

\begin{itemize}
  \item all \textbf{8 SoC-level assertions proved unboundedly}
        (2- and 3-inductive certificates), $\sim$7\,s each, $\approx$63\,s
        total, with memories unconstrained --- i.e.\ the invariants hold for
        \emph{every} program;
  \item \textbf{100\% COI coverage} of state bits by proven properties;
        12 mined conjectures were refuted by BMC and honestly reported as
        discarded;
  \item property authoring exposed \textbf{two real bugs} in the original
        SoC: (1)~\texttt{mkCsrNewVal} performs no WARL legalization, so
        \texttt{csrrw mstatus} can set \texttt{MPP}=2'b10 and a subsequent
        \texttt{mret} reaches the reserved privilege mode~2;
        (2)~no instruction-address-misaligned exception is raised and
        \texttt{mepc[1:0]} is not masked, so a misaligned PC is
        architecturally reachable while fetch silently truncates
        \texttt{fetchPC[1:0]}. Both were confirmed with BMC counterexamples
        and are documented in the benchmark's \texttt{PATCHES.md}.
\end{itemize}

\section{Evaluation}\label{sec:eval}

\paragraph{Setup.} 13 benchmarks (12 block-level designs spanning
arbiters, FIFOs, handshakes, Gray/LFSR counters, a 32-bit ALU, a memory
controller, a UART transmitter, an AXI-Lite slave, and a 3-stage forwarding
pipeline; plus the RV32IMA SoC), each with a validated-safe variant and 19
seeded-bug variants with \emph{measured} minimal counterexample depths
(2--38 cycles). Baseline: \tool{yosys-smtbmc} (the checking engine of
SymbiYosys) in BMC + temporal-induction portfolio, non-incremental Z3 (its
stronger configuration on this suite), on the \emph{identical} Yosys
elaboration. Both tools: Z3 4.x, same per-variant timeout, single core.
\sfv{} runs deterministically (\texttt{--no-llm}) for reproducibility; every
bug it reports is additionally replayed through Verilator.

\begin{table}[h]
\centering
\caption{Aggregate results: 32 variants (19 seeded bugs).
``Solved'' = expected verdict returned (unbounded proof for safe variants,
counterexample for bug variants).}
\label{tab:agg}
\begin{tabular}{lccc}
\toprule
Tool & Solved & Total wall time (s) & Bugs Verilator-confirmed \\
\midrule
\sfv{} (BMC + $k$-ind + Houdini) & \textbf{30/32} & 752 & 19/19 \\
\tool{yosys-smtbmc} (BMC + temporal ind.) & 30/32 & 530 & --- \\
\bottomrule
\end{tabular}
\end{table}


\begin{table}[h]
\centering
\caption{Per-variant results. State bits from the elaborated model;
``vlt'' = counterexample confirmed by Verilator replay.}
\label{tab:per}
\resizebox{\textwidth}{!}{%
\begin{tabular}{llccllrclr}
\toprule
Benchmark & Variant & Expect & Bits & \sfv{} & Method & t(s) & vlt &
\tool{smtbmc} & t(s) \\
\midrule
alu32 & design.sv & safe & 85 & \checkmark SAFE & kind & 0.1 &  & \checkmark SAFE & 19.0 \\
alu32 & bug1.sv & bug & 85 & \checkmark BUG & kind-base & 0.4 & \checkmark & \checkmark BUG & 0.2 \\
alu32 & bug2.sv & bug & 85 & \checkmark BUG & kind-base & 0.5 & \checkmark & \checkmark BUG & 0.1 \\
async\_handshake & design.sv & safe & 29 & \checkmark SAFE & kind & 0.8 &  & \checkmark SAFE & 1.7 \\
async\_handshake & bug1.sv & bug & 29 & \checkmark BUG & kind-base & 1.2 & \checkmark & \checkmark BUG & 0.2 \\
axi\_lite\_slave & design.sv & safe & 12 & \checkmark SAFE & kind & 0.1 &  & \checkmark SAFE & 1.3 \\
axi\_lite\_slave & bug1.sv & bug & 12 & \checkmark BUG & kind-base & 0.3 & \checkmark & \checkmark BUG & 0.1 \\
axi\_lite\_slave & bug2.sv & bug & 12 & \checkmark BUG & kind-base & 0.2 & \checkmark & \checkmark BUG & 0.1 \\
counter\_wrap & design.sv & safe & 10 & \checkmark SAFE & kind & 0.0 &  & \checkmark SAFE & 0.9 \\
counter\_wrap & bug1.sv & bug & 10 & \checkmark BUG & bmc & 1.4 & \checkmark & \checkmark BUG & 0.3 \\
gray\_counter & design.sv & safe & 29 & \checkmark SAFE & kind & 0.1 &  & \checkmark SAFE & 2.8 \\
gray\_counter & bug1.sv & bug & 29 & \checkmark BUG & bmc & 18.9 & \checkmark & \checkmark BUG & 1.6 \\
lfsr & design.sv & safe & 18 & \checkmark SAFE & kind & 0.0 &  & \checkmark SAFE & 1.1 \\
lfsr & bug1.sv & bug & 18 & \checkmark BUG & kind-base & 0.1 & \checkmark & \checkmark BUG & 0.1 \\
memctrl & design.sv & safe & 34 & -- UNKNOWN & bmc & 300.0 &  & -- UNKNOWN & 151.0 \\
memctrl & bug1.sv & bug & 34 & \checkmark BUG & kind-base & 0.2 & \checkmark & \checkmark BUG & 0.1 \\
memctrl & bug2.sv & bug & 34 & \checkmark BUG & kind-base & 0.2 & \checkmark & \checkmark BUG & 0.1 \\
pipeline\_hazard & design.sv & safe & 375 & -- UNKNOWN & bmc & 351.9 &  & -- UNKNOWN & 173.8 \\
pipeline\_hazard & bug1.sv & bug & 375 & \checkmark BUG & kind-base & 2.1 & \checkmark & \checkmark BUG & 0.2 \\
priority\_arbiter & design.sv & safe & 23 & \checkmark SAFE & kind & 0.7 &  & \checkmark SAFE & 3.3 \\
priority\_arbiter & bug1.sv & bug & 23 & \checkmark BUG & bmc & 11.9 & \checkmark & \checkmark BUG & 0.5 \\
round\_robin\_arbiter & design.sv & safe & 15 & \checkmark SAFE & kind & 0.1 &  & \checkmark SAFE & 2.4 \\
round\_robin\_arbiter & bug1.sv & bug & 15 & \checkmark BUG & kind-base & 0.1 & \checkmark & \checkmark BUG & 0.1 \\
round\_robin\_arbiter & bug2.sv & bug & 15 & \checkmark BUG & kind-base & 0.0 & \checkmark & \checkmark BUG & 0.1 \\
sync\_fifo & design.sv & safe & 15 & \checkmark SAFE & kind & 0.1 &  & \checkmark SAFE & 2.1 \\
sync\_fifo & bug1.sv & bug & 15 & \checkmark BUG & kind-base & 2.4 & \checkmark & \checkmark BUG & 0.5 \\
sync\_fifo & bug2.sv & bug & 15 & \checkmark BUG & kind-base & 0.1 & \checkmark & \checkmark BUG & 0.1 \\
uart\_tx & design.sv & safe & 28 & \checkmark SAFE & kind & 0.2 &  & \checkmark SAFE & 1.6 \\
uart\_tx & bug1.sv & bug & 28 & \checkmark BUG & bmc & 17.1 & \checkmark & \checkmark BUG & 2.1 \\
uart\_tx & bug2.sv & bug & 28 & \checkmark BUG & kind-base & 1.6 & \checkmark & \checkmark BUG & 0.2 \\
rv32i\_soc & design.sv & safe & 1878 & \checkmark SAFE & kind & 4.7 &  & \checkmark SAFE & 160.3 \\
rv32i\_soc & bug1.sv & bug & 1878 & \checkmark BUG & kind-base & 34.1 & \checkmark & \checkmark BUG & 1.4 \\
\bottomrule
\end{tabular}}
\end{table}


\paragraph{Scalability.} Table~\ref{tab:scal} orders the safe variants by model size.
\begin{table}[h]
\centering
\caption{Scalability: safe variants by state bits.}
\label{tab:scal}
\begin{tabular}{lrclcr}
\toprule
Benchmark & Bits & \sfv{} & t(s) & \tool{smtbmc} & t(s) \\
\midrule
counter\_wrap & 10 & SAFE & 0.0 & SAFE & 0.9 \\
axi\_lite\_slave & 12 & SAFE & 0.1 & SAFE & 1.3 \\
round\_robin\_arbiter & 15 & SAFE & 0.1 & SAFE & 2.4 \\
sync\_fifo & 15 & SAFE & 0.1 & SAFE & 2.1 \\
lfsr & 18 & SAFE & 0.0 & SAFE & 1.1 \\
priority\_arbiter & 23 & SAFE & 0.7 & SAFE & 3.3 \\
uart\_tx & 28 & SAFE & 0.2 & SAFE & 1.6 \\
async\_handshake & 29 & SAFE & 0.8 & SAFE & 1.7 \\
gray\_counter & 29 & SAFE & 0.1 & SAFE & 2.8 \\
memctrl & 34 & UNKNOWN & 300.0 & UNKNOWN & 151.0 \\
alu32 & 85 & SAFE & 0.1 & SAFE & 19.0 \\
pipeline\_hazard & 375 & UNKNOWN & 351.9 & UNKNOWN & 173.8 \\
rv32i\_soc & 1878 & SAFE & 4.7 & SAFE & 160.3 \\
\bottomrule
\end{tabular}
\end{table}


\paragraph{Discussion.} Both portfolios solve 30 of 32 variants: all 19
seeded bugs, and 11 of 13 safe variants. The two unsolved safe variants
(\texttt{memctrl}, \texttt{pipeline\_hazard}) defeat both engines'
unbounded methods for the same reason --- their inductive invariants require
quantified facts over memory contents and cross-stage datapath equalities
that neither temporal induction nor bound-style lemma templates express;
both tools honestly report a bounded result (proof radius $\geq 45$ resp.\
explicit timeout) rather than a claim. Beyond the tie in solve count, four
observations. (1)~\emph{Unbounded proofs scale differently}: on the SoC
safe variant \sfv{} proves all assertions in 4.7\,s against 160\,s for the
baseline ($34\times$), and 0.1--0.8\,s against 1--19\,s across the safe
blocks --- the engine finds the small inductive core once, while temporal
induction re-solves monolithically. (2)~\emph{Shallow-bug BMC favors the
baseline} (typically 0.1--2\,s vs.\ 0.1--34\,s): \tool{yosys-smtbmc}'s
compiled solver loop outpaces our Python-side incremental unrolling, and our
portfolio pays the induction attempt first; at the 300\,s budget this never
changes a verdict. (3)~\emph{Every} \sfv{} bug report --- 19 of 19 ---
replayed to a firing assertion under Verilator, including the memory
counterexamples once initial RAM state was aligned with two-state
simulation semantics (\S\ref{sec:cex}); the baseline offers no analogous
artifact. (4)~At SoC scale, word-level arrays keep the model compact; the
eight architectural invariants prove in seconds, indicating that the
practical bottleneck for full-SoC convergence is property authoring ---
precisely the step the harness automates.

\section{Related Work}

Commercial formal apps (Cadence JasperGold FPV, Synopsys VC~Formal, Siemens
Questa PropCheck) provide mature property proving but closed verdicts;
\sfv{} targets the same workflows with open, re-checkable artifacts.
SymbiYosys/\tool{yosys-smtbmc} is the standard open-source flow and serves
as our baseline. HWMCC defines the BTOR2 format we adopt. riscv-formal
pioneered reusable ISA-level property suites; our SoC benchmark follows its
spirit at the architectural-invariant level. Houdini originates in annotation
inference~\cite{houdini}; we combine it with LLM lemma proposal in the style
of recent AI-for-verification work, and with the artifact-first philosophy
of Pramaana Labs' everyday-provers program~\cite{pramaana}, which
demonstrated the same loop by formalizing tax law in Lean and surfacing a
marginal-relief bug with a proved fix.

\section{Limitations and Future Work}

The Lean artifacts are generated and pattern-matched against Sparkle's
actual API but not compiler-verified in this environment (no network route
to a Lean toolchain); wiring \texttt{lake build} into the agentic loop ---
with the LLM repairing elaboration errors --- is the natural next step, as
is discharging the emitted proof obligations with Lean tactics seeded by the
engine's invariants. The engine's unbounded power currently ends at Houdini
over template/LLM lemmas; IC3/PDR would extend it. Liveness is handled only
through safety encodings (bounded-wait counters). LLM-dependent stages were
exercised with the deterministic fallbacks in this evaluation because the
sandbox blocks the OpenRouter endpoint; the code paths are complete and
config-gated. Finally, mined-property precision depends on exploration
depth; refuted conjectures are reported, never silently dropped, but a
smarter sampler would raise the proposal hit-rate.

\section{Conclusion}

\sfv{} shows that the everyday-provers loop --- formalize, search with
provers, return inspectable artifacts, let AI accelerate but never vouch ---
is practical for hardware today: SystemVerilog in; Lean models, unbounded
certificates, Verilator-confirmed bugs, and a DV-grade convergence dossier
out. On a real Linux-booting SoC it proved the full architectural invariant
set in about a minute and, in the process of writing properties, found two
real bugs that a decade of simulation had not.

\begin{thebibliography}{9}
\bibitem{pramaana} Pramaana Labs. \emph{The Age of Everyday Provers}.
\url{https://pramaanalabs.ai/blog/the-age-of-everyday-provers}, 2026.
\bibitem{houdini} C.~Flanagan and K.~R.~M. Leino. Houdini, an annotation
assistant for ESC/Java. \emph{FME}, 2001.
\bibitem{btor2} A.~Niemetz, M.~Preiner, C.~Wolf, A.~Biere. BTOR2, BtorMC and
Boolector 3.0. \emph{CAV}, 2018.
\bibitem{yosys} C.~Wolf. Yosys Open SYnthesis Suite.
\url{https://yosyshq.net/yosys/}.
\bibitem{kind} M.~Sheeran, S.~Singh, G.~St{\aa}lmarck. Checking safety
properties using induction and a SAT-solver. \emph{FMCAD}, 2000.
\end{thebibliography}

\end{document}
